\chapter{Overview}
\label{chap:overview}

Let $\ps : \N \to \R_{\geq 0}$. The Duffin--Schaeffer conjecture, posed in 1941 and
proved by Koukoulopoulos and Maynard in 2019, says that if
\[ \sum_{q \geq 1} \frac{\ps(q)\varphi(q)}{q} = \infty \]
then almost every $\alpha \in [0,1]$ admits infinitely many coprime solutions to
\[ \left| \alpha - \frac{a}{q} \right| \leq \frac{\ps(q)}{q}. \tag{1.7} \]

This blueprint charts a formalization of that theorem, together with the
non-coprime dichotomy of the paper's Theorem 2. The terminal nodes are transcribed
verbatim from \texttt{ImperialCollegeLondon/AnnalsChallenge}, where they stand as
\texttt{sorry}; discharging them here discharges them there.

\section*{Why this target}

Mathlib has already done the first reduction and has left the gap explicitly marked.
\texttt{Mathlib/NumberTheory/WellApproximable.lean} (Oliver Nash, 2022) proves
\emph{Gallagher's ergodic theorem}: for any sequence of distances, the
well-approximable set is null or co-null. Its own module docstring reads:

\begin{quote}
Given a particular $\delta$, the Duffin--Schaeffer conjecture (now a theorem) gives a
criterion for deciding which of the two cases in the conclusion of Gallagher's theorem
actually occurs. \dots\ We do \emph{not} include a formalisation of the
Koukoulopoulos--Maynard result here.
\end{quote}

That is the hole this project fills. Because Gallagher gives a zero--one law, it
suffices to prove that the set has \emph{positive} measure --- and therefore no
constant anywhere in this blueprint ever has to be computed. That single observation
governs the shape of every estimate below.

\section*{State of Mathlib, September 2026}

Audited against the pinned revision. Present:

\begin{itemize}
  \item Gallagher's ergodic theorem, \texttt{AddCircle.addWellApproximable\_ae\_empty\_or\_univ}.
  \item Both Borel--Cantelli lemmas: \texttt{MeasureTheory.measure\_limsup\_cofinite\_eq\_zero}
    and \texttt{ProbabilityTheory.measure\_limsup\_eq\_one}. The second assumes genuine
    independence, which is exactly what fails here.
  \item Hausdorff measure $\mu H[d]$ and $\dimH$ with a usable API.
  \item \texttt{Nat.totient} with multiplicativity, \texttt{Nat.primesBelow}, Lebesgue
    density via Vitali families.
\end{itemize}

Absent, and needed:

\begin{itemize}
  \item The Chung--Erd\H{o}s inequality and the quasi-independent divergence
    Borel--Cantelli lemma. Searching the library for \texttt{chung}, \texttt{paley} or
    \texttt{second\_moment} returns nothing. Chapter \ref{chap:secondmoment}.
  \item Mertens' theorems. Searching for \texttt{mertens} returns only an unrelated
    Dedekind--Mertens TODO in a polynomial file. Chapter \ref{chap:anatomy}.
  \item The divisor bound and the average order of $\varphi$. Chapter \ref{chap:anatomy}.
  \item The mass transference principle of Beresnevich--Velani. Chapter
    \ref{chap:hausdorff}, and deliberately out of scope.
  \item Everything in Chapter \ref{chap:gcd}.
\end{itemize}

\section*{Scope}

Theorem 1 and Theorem 2(a),(b) are in scope. Corollary 3 --- the Hausdorff dimension
--- is not, and Chapter \ref{chap:hausdorff} says why: it needs mass transference,
which is an independent project of comparable size that would roughly double the work
while adding nothing to the conjecture itself.

\section*{How to read the status marks}

A node marked \notready\ has no Lean statement yet. Some of those are merely unwritten;
others are marked \textbf{TRANSCRIBE}, which means something stronger --- the statement
involves explicit exponents, normalisations or thresholds that must be copied from the
source rather than reconstructed. A plausible-looking Lean statement with a wrong
constant is worse than a visible gap, so those nodes carry their section reference and
nothing else until someone has the paper open.

\section*{Where to start}

Three independent entry points, in increasing order of commitment:

\begin{enumerate}
  \item Chapter \ref{chap:anatomy}. Mertens and the divisor bound depend on nothing
    else here and belong in Mathlib on their own merits.
  \item Chapter \ref{chap:secondmoment}. Chung--Erd\H{o}s is a Cauchy--Schwarz
    argument; it also belongs upstream.
  \item Theorem \ref{thm:two-a}, the convergence half of Theorem 2. It is the first
    Borel--Cantelli lemma applied to a majorant and needs none of Chapter
    \ref{chap:gcd}. It is the first node in this project whose proof is genuinely
    about Duffin--Schaeffer.
\end{enumerate}

Lemma \ref{lem:gallagher-no-tendsto} is a fourth: it is a TODO recorded in Mathlib's
own source, and closing it is a contribution to Mathlib independent of this project.

\chapter{The approximation sets}
\label{chap:setup}

\begin{definition}[Inequality (1.7)]
  \label{def:approximates}
  \lean{DuffinSchaeffer.Approximates}
  \leanok
  For $\ps : \N_{>0} \to \R_{\geq 0}$, $\alpha \in \R$, $a \in \N$ and
  $q \in \N_{>0}$, write $\mathrm{Approximates}(\ps, \alpha, a, q)$ for
  $\left| \alpha - a/q \right| \leq \ps(q)/q$.
\end{definition}

\begin{definition}[The sets $\Aq$]
  \label{def:setaq}
  \lean{DuffinSchaeffer.setAq}
  \leanok
  \uses{def:approximates}
  $\Aq = \{ \alpha \in [0,1] : \exists a,\ \gcd(a,q) = 1 \text{ and } (1.7) \text{ holds} \}$.
\end{definition}

\begin{definition}[$\A(\ps)$ and $\K(\ps)$]
  \label{def:seta}
  \lean{DuffinSchaeffer.setA, DuffinSchaeffer.setK}
  \leanok
  \uses{def:approximates}
  $\A(\ps)$ is the set of $\alpha \in [0,1]$ for which (1.7) has infinitely many
  coprime solutions $(a,q)$; $\K(\ps)$ is the same with coprimality dropped and
  $0 \leq a \leq q$ imposed.
\end{definition}

\begin{definition}[The radius $\Psi$ and the majorant $\pstar$]
  \label{def:psistar}
  \lean{DuffinSchaeffer.psiSup, DuffinSchaeffer.psiStar}
  \leanok
  $\Psi(q) = \sup \{ \ps(n)/n : q \mid n \}$ and $\pstar(q) = \varphi(q)\Psi(q)$.
  The supremum is taken in $[0,\infty]$, so it always exists.
\end{definition}

\begin{lemma}[$\Psi \leq \pstar$]
  \label{lem:psisup-le-psistar}
  \lean{DuffinSchaeffer.psiSup_le_psiStar}
  \leanok
  \uses{def:psistar}
  $\Psi(q) \leq \pstar(q)$, and hence $\sum_q \Psi(q) < \infty$ whenever
  $\sum_q \pstar(q) < \infty$.
\end{lemma}

\begin{proof}
  \leanok
  \textbf{Proved.} $\varphi(q) \geq 1$. Small, but it is what makes the covering argument
  of Theorem \ref{thm:two-a} work: the convergence hypothesis has to control the
  \emph{radii} $\Psi(d)$ and not only the weighted radii $\varphi(d)\Psi(d)$.
  Lean: \texttt{tsum\_psiSup\_lt\_top}.
\end{proof}

\begin{lemma}[$\Aq$ as a union of balls]
  \label{lem:setaq-union}
  \lean{DuffinSchaeffer.setAq_eq}
  \leanok
  \uses{def:setaq}
  $\Aq = [0,1] \cap \bigcup_{\gcd(a,q)=1} \overline{B}(a/q, \ps(q)/q)$.
\end{lemma}

\begin{proof}
  \leanok
  \textbf{Proved.} Unfolding both sides; (1.7) is membership of a closed ball.
\end{proof}

\begin{lemma}[Measurability of $\Aq$]
  \label{lem:measurable-setaq}
  \lean{DuffinSchaeffer.measurableSet_setAq}
  \leanok
  \uses{lem:setaq-union}
  $\Aq$ is measurable.
\end{lemma}

\begin{proof}
  \leanok
  \textbf{Proved.} A countable union of closed balls, intersected with $[0,1]$.
  No hypothesis on $\ps$ is needed.
\end{proof}

\begin{lemma}[Measure of a single $\Aq$]
  \label{lem:volume-setaq}
  \lean{DuffinSchaeffer.volume_setAq_le}
  \leanok
  \uses{lem:setaq-union, lem:truncate}
  If $\ps(q) \leq 1/2$ then $\Leb(\Aq) \leq 2\ps(q)\varphi(q)/q$.
\end{lemma}

\begin{proof}
  \leanok
  \textbf{Proved}, for the upper bound. Under $\ps(q) \leq 1/2$ no numerator outside
  $0 \leq a \leq q$ can bring a point of $[0,1]$ within $\ps(q)/q$ of $a/q$
  (\texttt{setAq\_subset}), so $\Aq$ is covered by $\#\{a \le q : \gcd(a,q)=1\}$ balls of
  radius $\ps(q)/q$; for $q \geq 2$ that count is exactly $\varphi(q)$, because $a = 0$
  and $a = q$ are both excluded by coprimality (\texttt{card\_numerators}).

  \emph{The hypothesis is not cosmetic, and this lemma was originally stated without it.}
  Two things go wrong. For large $\ps(q)$ the numerators $a > q$ also deposit mass in
  $[0,1]$, so the count $\varphi(q)$ is not an upper bound. And at $q = 1$ both $a = 0$
  and $a = 1$ are coprime to $1$ while $\varphi(1) = 1$, so $\#\{a \le q\} = \varphi(q)$
  simply fails; the conclusion still holds there, but only because the two intervals are
  truncated by $[0,1]$ to total length exactly $2\ps(1)$, which is a different argument
  (\texttt{volume\_setAq\_le\_one}). Both problems vanish under $\ps \leq 1/2$, which
  Lemma \ref{lem:truncate} supplies and everything downstream assumes anyway.

  The matching lower bound is now also proved, as
  \texttt{volume\_setAq\_of\_two\_le}: for $q \geq 2$ every ball lies inside $[0,1]$, so
  the containment above is an equality, and the balls are pairwise \emph{almost}
  disjoint --- at $\ps(q) = 1/2$ consecutive intervals share an endpoint, so the
  statement needed is a null intersection (\texttt{volume\_closedBall\_inter}) rather
  than disjointness, and \texttt{measure\_biUnion\_finset}$_0$ then gives
  $\Leb(\Aq) = 2\ps(q)\varphi(q)/q$ exactly.
\end{proof}

\begin{lemma}[Reduction to $\ps \leq 1/2$]
  \label{lem:truncate}
  \notready
  \uses{def:seta}
  Either $\Leb(\A(\ps)) = 1$ outright, or replacing $\ps$ by $\min(\ps, 1/2)$ changes
  nothing.
\end{lemma}

\begin{proof}
  If $\ps(q) > 1/2$ for infinitely many $q$, a single such $q$ already gives
  $\Aq = [0,1]$, so $\A(\ps) = [0,1]$. Otherwise the truncation alters only finitely
  many terms, which cannot affect a $\limsup$. Lean: \texttt{setA\_truncate}.
\end{proof}

\begin{lemma}[$\A$ as a $\limsup$ of the $\Aq$]
  \label{lem:seta-limsup}
  \lean{DuffinSchaeffer.mem_setA_iff}
  \leanok
  \uses{def:seta, def:setaq, lem:truncate}
  For $\ps \leq 1/2$, $\alpha \in \A(\ps)$ iff $\alpha \in [0,1]$ and $\alpha$ lies in
  infinitely many $\Aq$.
\end{lemma}

\begin{proof}
  The content is that infinitely many \emph{pairs} $(a,q)$ forces infinitely many
  \emph{denominators}. Under $\ps \leq 1/2$ a solution has $0 \leq a \leq q$, so each
  denominator admits only finitely many numerators. Without that hypothesis the
  implication genuinely fails, which is why Lemma \ref{lem:truncate} comes first.
  Lean: \texttt{mem\_setA\_iff}.
\end{proof}

\begin{lemma}[Measurability of $\A$ and $\K$]
  \label{lem:measurable-seta}
  \lean{DuffinSchaeffer.measurableSet_setA, DuffinSchaeffer.measurableSet_setK}
  \leanok
  \uses{def:seta}
  $\A(\ps)$ and $\K(\ps)$ are measurable.
\end{lemma}

\begin{proof}
  \leanok
  \textbf{Proved.} Both are $[0,1]$ intersected with a set of the form
  $\{\alpha : \{i : \alpha \in B_i\} \text{ infinite}\}$ for a countable family --- here
  indexed by pairs $(a,q)$, each $B_{(a,q)}$ a closed ball or empty. Such a set equals
  $\bigcap_{t} \bigcup_{i \notin t} B_i$ over the finite subsets $t$, which is a countable
  intersection of countable unions because $\mathrm{Finset}\,\iota$ is countable when
  $\iota$ is. Stated once as \texttt{measurableSet\_setOf\_infinite} and used for both.
  No hypothesis on $\ps$ is needed.
\end{proof}

\begin{lemma}[Divergence transfers to the measures]
  \label{lem:divergence}
  \lean{DuffinSchaeffer.not_summable_volume_setAq}
  \leanok
  \uses{lem:volume-setaq, lem:truncate}
  If $\sum \ps(q)\varphi(q)/q$ diverges and $\ps \leq 1/2$, then
  $\sum_q \Leb(\Aq) = \infty$.
\end{lemma}

\begin{proof}
  \leanok
  \textbf{Proved.} Immediate from the equality in Lemma \ref{lem:volume-setaq} once the
  $q = 1$ term is dealt with. It is absorbed rather than computed: the equality is
  available for $q \geq 2$, and a single finite term cannot rescue a divergent series,
  so bounding $\Leb(\Aq)$ below by the $q \ge 2$ formula plus an indicator at $q = 1$
  suffices.
\end{proof}

\chapter{The zero--one law}
\label{chap:gallagher}

\begin{theorem}[Gallagher's ergodic theorem]
  \label{thm:gallagher}
  \lean{AddCircle.addWellApproximable_ae_empty_or_univ}
  \leanok
  \mathlibok
  For any $\delta : \N \to \R$ tending to $0$, the set
  $\mathrm{addWellApproximable}(\R/\Z, \delta)$ is either null or co-null.
\end{theorem}

\begin{proof}
  \leanok
  In Mathlib, \texttt{Mathlib/NumberTheory/WellApproximable.lean}. Nothing to do.
\end{proof}

\begin{definition}[The distance sequence]
  \label{def:distseq}
  \lean{DuffinSchaeffer.distSeq}
  \leanok
  $\delta_\ps(n) = \ps(n)/n$ for $n > 0$, and $0$ at $n = 0$.
\end{definition}

\begin{lemma}[Dropping the vanishing hypothesis]
  \label{lem:gallagher-no-tendsto}
  \notready
  \uses{thm:gallagher}
  Gallagher's theorem holds for any non-negative $\delta$, with no assumption that
  $\delta \to 0$.
\end{lemma}

\begin{proof}
  Mathlib's own TODO records the argument: if $\delta \not\to 0$ then
  $\mathrm{addWellApproximable}(\R/\Z,\delta)$ is all of $\R/\Z$, by an elementary
  non-measure-theoretic argument. The Duffin--Schaeffer hypothesis does not supply
  $\delta \to 0$, so this cannot be skipped. It is a self-contained Mathlib
  contribution. Lean: \texttt{addWellApproximable\_ae\_empty\_or\_univ\_of\_nonneg}.
\end{proof}

\begin{lemma}[Transfer to the circle]
  \label{lem:transfer}
  \notready
  \uses{def:seta, def:distseq}
  For $\ps \leq 1/2$, $\A(\ps)$ agrees up to a null set with the preimage of
  $\mathrm{addWellApproximable}(\R/\Z, \delta_\ps)$ under $\R \to \R/\Z$,
  intersected with $[0,1]$.
\end{lemma}

\begin{proof}
  Bookkeeping: the quotient is measure-preserving on $[0,1)$, and
  \texttt{UnitAddCircle.mem\_addWellApproximable\_iff} already expresses membership in
  the concrete form $\exists m < n,\ \gcd(m,n) = 1,\ \|x - m/n\| < \delta_n$. The
  discrepancy between $<$ and $\leq$ costs a null set. Lean:
  \texttt{setA\_eq\_preimage\_addWellApproximable}.
\end{proof}

\begin{proposition}[Positive measure suffices]
  \label{prop:positive-suffices}
  \notready
  \uses{thm:gallagher, lem:gallagher-no-tendsto, lem:transfer}
  If $\Leb(\A(\ps)) > 0$ then $\Leb(\A(\ps)) = 1$.
\end{proposition}

\begin{proof}
  Immediate from Lemmas \ref{lem:transfer} and \ref{lem:gallagher-no-tendsto}. This is
  the proposition that makes every constant in the rest of the blueprint irrelevant.
  Lean: \texttt{volume\_setA\_eq\_one\_of\_pos}.
\end{proof}

\chapter{The second moment method}
\label{chap:secondmoment}

\begin{lemma}[Chung--Erd\H{o}s]
  \label{lem:chung-erdos}
  \notready
  For finitely many measurable events,
  $\left(\sum_i \mu(A_i)\right)^2 \leq \mu\!\left(\bigcup_i A_i\right) \sum_{i,j} \mu(A_i \cap A_j)$.
\end{lemma}

\begin{proof}
  Cauchy--Schwarz applied to $\mathbf{1}_{\bigcup A_i}$ and $\sum_i \mathbf{1}_{A_i}$.
  Not in Mathlib; belongs upstream. Lean: \texttt{chung\_erdos}.
\end{proof}

\begin{proposition}[Divergence Borel--Cantelli, quasi-independent form]
  \label{prop:quasi-bc}
  \notready
  \uses{lem:chung-erdos}
  Let $(A_i)$ be measurable, and suppose there is a constant $C$ and a sequence of
  finite index sets $S_n$ with $\sum_{i \in S_n} \mu(A_i) \to \infty$ and
  \[ \sum_{i,j \in S_n} \mu(A_i \cap A_j) \leq C \left( \sum_{i \in S_n} \mu(A_i) \right)^2 \]
  for every $n$. Then the set of points lying in infinitely many $A_i$ has measure at
  least $C^{-1}$.
\end{proposition}

\begin{proof}
  Chung--Erd\H{o}s on $S_n$ gives $\mu(\bigcup_{i \in S_n} A_i) \geq C^{-1}$; the tail
  sets inherit the bound because the hypothesis is stable under discarding an initial
  segment, and $\mu$ is continuous from above on a finite measure space.

  The quantifier structure matters and is the reason this is stated with a
  \emph{sequence} of finite sets. Quasi-independence over all finite sets of
  denominators is false, and its failure is precisely the difficulty of the
  conjecture. Lean: \texttt{measure\_infinite\_pos\_of\_overlap}.
\end{proof}

\chapter{Anatomy of integers}
\label{chap:anatomy}

Nothing in this chapter is due to Koukoulopoulos and Maynard; it is the standard
toolkit their argument assumes, and Mathlib does not have it. Every node is
independently upstreamable and none depends on the rest of the project.

\begin{lemma}[Mertens' second theorem]
  \label{lem:mertens-two}
  \notready
  $\sum_{p < n} 1/p = \log\log n + M + O(1/\log n)$.
  Lean: \texttt{mertens\_second}.
\end{lemma}

\begin{lemma}[Mertens' third theorem]
  \label{lem:mertens-three}
  \notready
  \uses{lem:mertens-two}
  $\prod_{p < n} (1 - 1/p)^{-1} \asymp \log n$. The sharp constant $e^{\gamma}$ is
  never needed. Lean: \texttt{mertens\_third}.
\end{lemma}

\begin{lemma}[Divisor bound]
  \label{lem:divisor-bound}
  \notready
  $d(n) = O_{\varepsilon}(n^{\varepsilon})$. Lean: \texttt{divisor\_bound}.
\end{lemma}

\begin{lemma}[Dirichlet swap]
  \label{lem:dirichlet-swap}
  \lean{DuffinSchaeffer.sum_sum_divisors}
  \leanok
  For any $f$, $\sum_{q \leq n} \sum_{d \mid q} f(d) = \sum_{d \leq n} \lfloor n/d \rfloor f(d)$.
\end{lemma}

\begin{proof}
  \leanok
  \textbf{Proved.} The divisors of $q$ are exactly the $d \leq n$ dividing $q$ once
  $0 < q \leq n$ (\texttt{divisors\_eq\_filter}), so both sides count the same pairs;
  exchanging the order of summation and counting multiples with
  \texttt{Nat.Ioc\_filter\_dvd\_card\_eq\_div} gives the result.
\end{proof}

\begin{lemma}[Average order of $\varphi$]
  \label{lem:totient-average}
  \lean{DuffinSchaeffer.sum_totient_div_self}
  \leanok
  \uses{lem:dirichlet-swap}
  $\tfrac{n}{2} \leq \sum_{q \leq n} \varphi(q)/q \leq n$; in particular
  $\sum_{q \leq n} \varphi(q)/q \asymp n$.
\end{lemma}

\begin{proof}
  \leanok
  \textbf{Proved}, with explicit constants $c = 1/2$ and $C = 1$.

  The upper bound is $\varphi(q) \leq q$ termwise. For the lower bound, sum
  $\sum_{d \mid q} \varphi(d) = q$ (Mathlib's \texttt{Nat.sum\_totient}) over $q \leq n$.
  The left side is $n(n+1)/2$ by Gauss; Lemma \ref{lem:dirichlet-swap} rewrites it as
  $\sum_{d \leq n} \lfloor n/d \rfloor \varphi(d)$, which is at most
  $n \sum_{d \leq n} \varphi(d)/d$. Dividing by $n$ gives
  $\sum_{q \leq n} \varphi(q)/q \geq (n+1)/2$.

  Worth recording because the obvious route is much worse: expanding
  $\varphi(q)/q = \sum_{d \mid q} \mu(d)/d$ leads to $\sum_{d} \mu(d)/d^2 = 1/\zeta(2)$
  and the Basel problem, none of which is needed. The constant $1/2$ is not sharp --- the
  truth is $6/\pi^2 \approx 0.61$ --- but no estimate downstream cares.
\end{proof}

\chapter{Overlap estimates}
\label{chap:overlap}

\begin{lemma}[Pollington--Vaughan]
  \label{lem:pollington-vaughan}
  \notready
  \uses{lem:volume-setaq, lem:mertens-three}
  \textbf{TRANSCRIBE} --- Pollington and Vaughan, \emph{Mathematika} 37 (1990),
  190--200; quoted in the paper's \S2.

  For $q \neq r$ with $d = \gcd(q,r)$, the overlap $\Leb(\Aq \cap A_r)$ exceeds the
  independent prediction $\Leb(\Aq)\Leb(A_r)$ by a factor controlled by a product
  $\prod (1-1/p)^{-1}$ over the primes $p$ dividing $qr/d^2$ that exceed a threshold
  depending on $q$, $r$ and $\ps$.

  The threshold carries normalisations that must be copied from the source. It is
  deliberately not stated in Lean until then: with the threshold taken too large the
  statement is false, with it taken too small the statement is useless, and neither
  error is visible downstream until the whole chapter has been built on it.
\end{lemma}

\begin{theorem}[Selection of quasi-independent denominators]
  \label{thm:overlap-core}
  \notready
  \uses{lem:pollington-vaughan, thm:gcd-iteration, lem:divergence}
  Suppose $\ps \leq 1/2$ and $\sum_q \Leb(\Aq) = \infty$. Then there are a constant
  $C$ and finite sets $S_n$ of denominators with
  $\sum_{q \in S_n} \Leb(\Aq) \to \infty$ and
  \[ \sum_{q,r \in S_n} \Leb(\Aq \cap A_r) \leq C \left( \sum_{q \in S_n} \Leb(\Aq) \right)^2 . \]
\end{theorem}

\begin{proof}
  This is the entire content of the paper from \S2 onwards, and the only node in the
  project with no precedent anywhere. Lemma \ref{lem:pollington-vaughan} converts the
  overlaps into a question about the small prime factors of $qr/\gcd(q,r)^2$; Theorem
  \ref{thm:gcd-iteration} answers that question.

  The constant $C$ is never computed: Proposition \ref{prop:quasi-bc} turns it into a
  positive lower bound $C^{-1}$ on the measure, and Proposition
  \ref{prop:positive-suffices} then upgrades that to $1$.
  Lean: \texttt{exists\_quasi\_independent\_subsets}.
\end{proof}

\chapter{GCD graphs}
\label{chap:gcd}

The machinery of the paper, \S\S3--10. The overlap of $\Aq$ and $A_r$ is inflated by
the small primes dividing $qr/\gcd(q,r)^2$, so bounding the double sum means
controlling how often a set of denominators can be pairwise entangled in that sense.
Koukoulopoulos and Maynard encode the entanglement as a weighted graph, attach a
numerical \emph{quality} to it, and show that a graph which is bad for the overlap
estimate admits a modification strictly increasing its quality. Quality is bounded
above, so the iteration terminates, and the terminal graphs are exactly those for
which the estimate holds.

\begin{definition}[GCD graph]
  \label{def:gcdgraph}
  \lean{DuffinSchaeffer.GCDGraph}
  \leanok
  A finite vertex set $V$ of denominators, a symmetric irreflexive edge set $E$ on it,
  a finite set $P$ of primes, and for each $p \in P$ exponents $f_p, g_p$ with
  $g_p \leq f_p$ such that $p^{f_p} \| v$ for every $v \in V$ and
  $p^{g_p} \| \gcd(v,w)$ for every edge $(v,w)$.
\end{definition}

\begin{definition}[Quality]
  \label{def:quality}
  \notready
  \uses{def:gcdgraph}
  \textbf{TRANSCRIBE} --- paper \S3.

  A numerical invariant $\quality(\gcdgraph)$ of a GCD graph, built from the edge
  density and the exponents $f_p, g_p$. The explicit normalisation must be copied from
  the paper; the entire argument is an optimisation of this quantity, so a
  reconstructed formula would be undetectably wrong.
\end{definition}

\begin{lemma}[Quality is bounded]
  \label{lem:quality-bounded}
  \notready
  \uses{def:quality}
  \textbf{TRANSCRIBE} --- paper \S3. This is what makes the iteration terminate.
\end{lemma}

\begin{lemma}[Graph operations]
  \label{lem:gcd-operations}
  \notready
  \uses{def:quality}
  \textbf{TRANSCRIBE} --- paper \S\S5--9. The catalogue of modifications, each
  shown to increase quality: passing to an induced subgraph on a large vertex subset,
  enlarging $P$ and refining $f$ and $g$ along a new prime, and the compression step
  that removes vertices of unbalanced degree.
\end{lemma}

\begin{theorem}[The iteration]
  \label{thm:gcd-iteration}
  \notready
  \uses{lem:quality-bounded, lem:gcd-operations, lem:divisor-bound, lem:mertens-two}
  \textbf{TRANSCRIBE} --- paper \S10. A GCD graph of maximal quality has the structure
  needed for the overlap estimate; iterating the operations of Lemma
  \ref{lem:gcd-operations} from an arbitrary graph reaches one.
\end{theorem}

\chapter{Assembly}
\label{chap:main}

\begin{theorem}[Duffin--Schaeffer]
  \label{thm:one}
  \notready
  \uses{prop:positive-suffices, prop:quasi-bc, thm:overlap-core, lem:seta-limsup}
  If $\sum \ps(q)\varphi(q)/q$ diverges then $\A(\ps)$ is measurable and
  $\Leb(\A(\ps)) = 1$.
\end{theorem}

\begin{proof}
  Truncate to $\ps \leq 1/2$ by Lemma \ref{lem:truncate}. Lemma \ref{lem:divergence}
  gives $\sum \Leb(\Aq) = \infty$; Theorem \ref{thm:overlap-core} supplies the sets
  $S_n$ and the constant $C$; Proposition \ref{prop:quasi-bc} gives
  $\Leb(\A(\ps)) \geq C^{-1} > 0$ via Lemma \ref{lem:seta-limsup}; Proposition
  \ref{prop:positive-suffices} upgrades it to $1$. Lean: \texttt{theorem\_1}.
\end{proof}

\begin{theorem}[Theorem 2(a), convergence]
  \label{thm:two-a}
  \lean{DuffinSchaeffer.theorem_2_a}
  \leanok
  \uses{def:psistar, def:seta}
  If $\sum_q \pstar(q) < \infty$ then $\K(\ps)$ is measurable and $\Leb(\K(\ps)) = 0$.
\end{theorem}

\begin{proof}
  \leanok
  \uses{lem:measurable-seta, lem:psisup-le-psistar}
  \textbf{Proved.} The first terminal node of the project to close. What follows is the
  argument as formalised.
  Measurability is done (Lemma \ref{lem:measurable-seta}). For the null statement:

  Given $\alpha \in \K(\ps)$ there are infinitely many pairs $(a,q)$, $a \leq q$, with
  $|\alpha - a/q| \leq \ps(q)/q$. Reduce $a/q$ to lowest terms $b/d$, $d \mid q$; then
  $|\alpha - b/d| \leq \ps(q)/q \leq \Psi(d)$. So $\alpha$ lies in
  \[ E_d := \bigcup_{\substack{0 \le b \le d \\ \gcd(b,d)=1}} \overline{B}(b/d, \Psi(d)), \qquad
     \Leb(E_d) \leq 2(\varphi(d)+1)\Psi(d) \leq 4\pstar(d), \]
  the last step by Lemma \ref{lem:psisup-le-psistar}. Since $\sum_d \pstar(d) < \infty$,
  the first Borel--Cantelli lemma (Mathlib:
  \texttt{MeasureTheory.measure\_limsup\_cofinite\_eq\_zero}) kills $\limsup_d E_d$.

  \textbf{The gap.} That argument needs infinitely many \emph{distinct} $d$, and the
  reduction to lowest terms does not supply them: infinitely many pairs can reduce to a
  single $b/d$. Take $\alpha = 0$, where every pair $(0,q)$ reduces to $0/1$. So the
  step ``infinitely many pairs, hence infinitely many $d$'' is false as stated.

  Two repairs are needed, and both are cheap once seen. First, discard the rationals:
  $\mathbb{Q}$ is countable, hence null, so it suffices to treat irrational $\alpha$.
  Second, note that $\sum_q \pstar(q) < \infty$ forces $\pstar(q) \to 0$, and
  $\pstar(q) \geq \Psi(q) \geq \ps(q)/q$, so $\ps(q)/q \to 0$. For irrational $\alpha$
  and fixed $b/d$ we have $|\alpha - b/d| > 0$, so only finitely many $q$ can satisfy
  $\ps(q)/q \geq |\alpha - b/d|$ --- that is, each $(b,d)$ absorbs only finitely many
  pairs. Infinitely many pairs therefore do give infinitely many $(b,d)$, and since
  $b \leq d$ bounds the numerators, infinitely many $d$.

  Worth stating plainly because the flawed version is the one that comes to mind first
  and it is not visibly wrong.

  Both repairs are in the Lean proof: \texttt{finite\_ratio\_ge} is the statement that
  $\{q : \delta \le \ps(q)/q\}$ is finite for $\delta > 0$, obtained from
  \texttt{ENNReal.finite\_const\_le\_of\_tsum\_ne\_top}, and the rationals are removed by
  \texttt{Set.Countable.measure\_zero}. The reduction to lowest terms is
  \texttt{redNum}/\texttt{redDen}, and the final step is Mathlib's
  \texttt{measure\_limsup\_cofinite\_eq\_zero} composed with
  \texttt{Filter.cofinite.limsup\_set\_eq}.
\end{proof}

\begin{theorem}[Theorem 2(b), divergence]
  \label{thm:two-b}
  \notready
  \uses{thm:one, def:psistar}
  If $\sum_q \pstar(q) = \infty$ then $\Leb(\K(\ps)) = 1$.
\end{theorem}

\begin{proof}
  Reduce to Theorem \ref{thm:one} by applying it to a function manufactured from
  $\pstar$, chosen so that its coprime approximation set embeds in $\K(\ps)$.
  Lean: \texttt{theorem\_2\_b}.
\end{proof}

\chapter{Out of scope: Hausdorff dimension}
\label{chap:hausdorff}

\begin{definition}[Critical exponent]
  \label{def:critexp}
  \lean{DuffinSchaeffer.critExp}
  \leanok
  $s = \inf \{ \beta : \sum_q \varphi(q)(\ps(q)/q)^{\beta} < \infty \}$.
\end{definition}

\begin{corollary}[Corollary 3 of the paper]
  \label{cor:three}
  \notready
  \uses{thm:one, def:critexp}
  For $\ps$ taking values in $[0,1/2]$, $\dimH \A(\ps) = \min(s,1)$.
  Lean: \texttt{corollary\_3}.
\end{corollary}

This does not follow from Theorem \ref{thm:one} by soft arguments. The passage from
full Lebesgue measure to full Hausdorff dimension runs through the mass transference
principle of Beresnevich and Velani (\emph{Annals} 164 (2006), 971--992), which
Mathlib does not have --- searching for \texttt{massTransference} or
\texttt{beresnevich} returns nothing --- and which brings its own prerequisite stack:
Hausdorff content, a Cantor-set construction, the ubiquity framework.

That is a separate project of comparable size which happens to share a statement file.
Folding it into the critical path would roughly double the work and add nothing to the
conjecture. Mathlib does have $\dimH$ and $\mu H[d]$ with a workable API, so the
landing site exists whenever someone wants it.
